% Evans 9.4: printed pp.514--521; English PDF pp.531--538.
% Chinese: printed p.363 / scan p.381.
\section{Nonexistence of Solutions}
\label{sec:9-4}
\Needspace{40mm}
\subsection{Large data: a weighted moment becomes infinite}
Assume $U$ is connected and let $\phi_1>0$ be the first Dirichlet
eigenfunction, $-\Delta\phi_1=\lambda_1\phi_1$, normalized by
$\int_U\phi_1=1$. For the problem $u_t-\Delta u=u^2$,
$u|_{\partial U}=0$, $u(0)=g\geq0$, take a classical solution with enough
boundary regularity for Green's identity. The maximum principle gives
$u\geq0$ as long as it exists.
\Needspace{50mm}
\begin{theorem}[Large-data obstruction; source \S9.4.1, Theorem 1]
If $\eta_0=\int g\phi_1>\lambda_1$, a classical solution cannot persist
through the time
\[
 T_*={1\over\lambda_1}\log{\eta_0\over\eta_0-\lambda_1}.
\]
This is an upper bound for its lifespan, not necessarily its actual blow-up time.
\end{theorem}
\begin{expandedproof}[normalization gives Jensen's inequality]
Set $\eta(t)=\int u(t)\phi_1$. Both $u$ and $\phi_1$ vanish on the
boundary, so the two boundary terms in Green's second identity vanish.
Consequently
\[
 \eta'=-\lambda_1\eta+\int u^2\phi_1
 \geq-\lambda_1\eta+\eta^2.
\]
The last step uses the probability measure $\phi_1dx$. Comparison with
the scalar ODE, or integration of $\xi=e^{\lambda_1t}\eta$, gives
\[
 \eta(t)\geq {\lambda_1\eta_0e^{-\lambda_1t}
 \over\lambda_1-\eta_0(1-e^{-\lambda_1t})}\qquad(0\leq t<T_*).
\]
The denominator stays positive until $T_*$ and tends to zero there.
If the PDE solution existed smoothly through that time, its finite moment
would contradict this bound. It might cease to exist earlier. Moreover
$\eta(t)\leq\|u(t)\|_\infty$, so moment divergence implies amplitude
divergence when the solution persists up to the comparison time.
\end{expandedproof}
\PDEMarginDiagram{9.M}{反应与扩散的阈值}
{归一化为 $\lambda_1=1$ 后，比较 ODE 为 $y'=y(y-1)$。箭头只表示这个 ODE 的相线；$y_0\leq1$ 时不能据此断言 PDE 全局存在。}
{ch09-blowup-margin}

\Needspace{40mm}
\subsection{Subcritical Fujita blow-up in all of space}
\Needspace{50mm}
\begin{theorem}[Small-data obstruction; source Theorem 2]
Let $g\in C_c^\infty(\R^n)$ be nonnegative and nonzero. If
$1<p<1+2/n$, there is no global nonnegative classical solution of
$u_t-\Delta u=u^p$ in the class $C([0,T];L^1\cap L^\infty)$ for every
finite $T$. This class supplies the decay and integrability for the calculation.
\end{theorem}
\begin{expandedproof}[a scale selected after the initial datum]
Fix $s>0$, put $\Phi_s(x)=(4\pi s)^{-n/2}e^{-|x|^2/(4s)}$, and set
$\eta(t)=\int u(t)\Phi_s$. Since
\[
 \Delta\Phi_s=\left({|x|^2\over4s^2}-{n\over2s}\right)\Phi_s,
 \qquad \int\Phi_s=1,
\]
Green's identity and Jensen give
$\eta'\geq-\lambda\eta+\eta^p$ with $\lambda=n/(2s)$.
To justify the identity on $\R^n$, test first with $\chi_R\Phi_s$.
After integrating by parts onto this compactly supported test, the errors
$u(2D\chi_R\cdot D\Phi_s+\Phi_s\Delta\chi_R)$ tend to zero.
On each finite time interval, the assumed $L^1,L^\infty$ bounds also
control $u^p$; dominated convergence gives the integral differential
identity. Thus no unproved decay of $Du$ is required.

The comparison ODE blows up if $\eta(0)^{p-1}>\lambda$.
This condition is equivalent to
\[
 (4\pi)^{-n/2}\int g(x)e^{-|x|^2/(4s)}dx
 > (n/2)^{1/(p-1)}s^{\,n/2-1/(p-1)}.
\]
As $s\to\infty$, the left side tends to a positive constant and the right
side tends to zero, since $p<1+2/n$. Choose one such finite $s$; the
ODE comparison then forbids a global solution. At $p=1+2/n$ the right
side is constant: this argument does not establish the endpoint.
\end{expandedproof}
\Needspace{42mm}
\PDEAnnotation{两个临界指数}{热方程这里出现 $1+2/n$，来自质量与扩散尺度；椭圆 Pohozaev 中出现 $(n+2)/(n-2)$，来自 Dirichlet 能量与幂次能量的尺度。两者不能混同。}

\Needspace{40mm}
\subsection{The full Derrick--Pohozaev identity}
Let $n\geq3$, $f\in C(\R)$, $F(s)=\int_0^sf(r)dr$, and
$u\in C^2(\overline U)$ solve $-\Delta u=f(u)$ with zero boundary data.
\Needspace{50mm}
\begin{theorem}[Identity and nonexistence; source \S9.4.2, Theorem 3]
One has
\begin{equation}\label{eq:9-pohozaev}
 {n-2\over2}\int_U|Du|^2
 +{1\over2}\int_{\partial U}(x\cdot\nu)(\partial_\nu u)^2
 =n\int_U F(u).
\end{equation}
If $U$ is star-shaped about $0$ and $f(u)=|u|^{p-1}u$ with
$p>(n+2)/(n-2)$, then $u\equiv0$.
\end{theorem}
\begin{expandedproof}[every boundary term retained]
Introduce the vector field
$Q=(x\cdot Du)Du-x|Du|^2/2$. Direct differentiation cancels the Hessian
cross terms and yields
\[
 \operatorname{div}Q=(x\cdot Du)\Delta u+{2-n\over2}|Du|^2.
\]
On $\partial U$, all tangential derivatives of $u$ vanish, so
$Du=(\partial_\nu u)\nu$ and
$Q\cdot\nu=(x\cdot\nu)(\partial_\nu u)^2/2$.
Also $f(u)x\cdot Du=x\cdot DF(u)$, and
$\int x\cdot DF(u)=-n\int F(u)$ since $F(0)=0$ on the boundary.
The divergence theorem gives \eqref{eq:9-pohozaev} with the displayed signs.

For a $C^1$ star-shaped boundary, choose a local defining function
$\rho$ with $U=\{\rho<0\}$. Since $(1-t)x\in\overline U$ for small
$t>0$, $\rho((1-t)x)\leq0=\rho(x)$, so
$-D\rho(x)\cdot x\leq0$; hence $x\cdot\nu\geq0$.
Testing the PDE by $u$ gives
$\int|Du|^2=\int|u|^{p+1}$. Substitution in the identity implies
\[
 \left({n-2\over2}-{n\over p+1}\right)\int|u|^{p+1}\leq0.
\]
The coefficient is strictly positive at a supercritical exponent, so
$u=0$. At equality of the exponents it is zero: the displayed inequality
alone gives no conclusion. On a strictly star-shaped domain a positive
nonzero solution at the critical exponent is also ruled out by the
boundary term and Hopf's lemma, but this is a separate argument and does
not extend the stated theorem to arbitrary sign-changing critical solutions.
\end{expandedproof}
\Needspace{42mm}
\PDEAnnotation{边界几何进入哪一步}{星形性只用来保证 $x\cdot\nu\geq0$。Pohozaev 恒等式本身对一般光滑区域成立；若区域不是星形，不能直接丢掉边界积分。}
